\section{Results}

Every number in this section is generated from a file in \texttt{results/} by
\texttt{benchmarks/paper\_numbers.py}; none is typed into the text. Figures are
rendered from the same files.

\subsection{Component validation}

Each component is checked against an external reference before it is used for
anything. Descriptors are compared against reference columns distributed with the
benchmark datasets, over \descriptorMolecules{} molecules, with the best agreement
reaching \descriptorBest{} exact. Free-energy estimators are run on an analytically
solvable switching problem where $\Delta F$ is known in closed form. Persistent
homology is run on point clouds of known topology --- a circle must yield exactly one
long-lived $H_1$ feature, acyclic octane none. Every autodiff operation is checked
against central finite differences. The Gene Ontology layer loads \goTerms{} terms
and \goAnnotations{} annotations, recovers known pathway enrichments as positive
controls, and on \goFpTrials{} random gene sets of \goFpGenes{} genes returns a mean
of \goFpMeanTerms{} significant terms, with any hit at all in \goFpAnyHit{} of trials
--- the false-positive behaviour the Benjamini--Hochberg procedure
\cite{benjamini1995controlling} promises. The enrichment the Mantis panel runs is a
separate, pruned index and is validated separately (Section~\ref{sec:go}).

\subsection{Property prediction}
\label{sec:molnet}

\begin{figure}[t]
\centering
\includegraphics[width=\linewidth]{fig1_molnet_benchmark}
\caption{MoleculeNet \cite{wu2018moleculenet} under scaffold-balanced splits
\cite{yang2019analyzing}, three seeds, with seed-to-seed spread. The neural model is
competitive but not dominant.}
\label{fig:molnet}
\end{figure}

The D-MPNN reproduces published numbers under the same protocol: BACE
\molnetBaceDMpnn{} against a published 0.857, BBBP \molnetBbbpDMpnn{} against
0.913, ClinTox \molnetClintoxDMpnn{} against 0.897
\cite{yang2019analyzing}. That is the check that the implementation is right.

It is not the headline. On BACE a plain ECFP-plus-descriptor random forest reaches
\molnetBaceEcfpDescriptorsRf{}, and on BBBP an ECFP random forest reaches
\molnetBbbpEcfpRf{} --- both ahead of the message-passing network, on datasets of one
to two thousand molecules.

Where the neural model wins is where the regime argument says it should. On
Lipophilicity, the largest set here at 4\,200 molecules, the D-MPNN reaches
\molnetLipoDMpnn{} RMSE: ahead of every fingerprint baseline
(\molnetLipoEcfpDescriptorsRf{} for the best of them) and ahead of the published
0.655 for the same architecture and protocol \cite{yang2019analyzing}. It also
wins on ClinTox (\molnetClintoxDMpnn{} against \molnetClintoxEcfpRf{}). A learned
representation needs data to pay for itself, and across these six datasets the
crossover is visible rather than asserted.

Both halves are reported because quoting only the dataset where the fancier model
wins is how benchmark tables become untrustworthy.

\subsection{Binding affinity}

On the Davis kinase panel \cite{davis2011comprehensive}, \affinityPairs{} measured
pairs over \affinityDrugs{} drugs and \affinityTargets{} targets, a random split
gives concordance index \affinityRandomCi{} and Pearson \affinityRandomPearson{},
against roughly 0.878 for DeepDTA on the same dataset \cite{ozturk2018deepdta}.

The number that matters for prospective use is the other one. Held out by
\emph{drug} --- the realistic case, where the compound has never been assayed against
anything --- concordance falls to \affinityColddrugCi{} and Pearson to
\affinityColddrugPearson{}. Held out by target it falls less far, to
\affinityColdtargetCi{}. A random split on an interaction matrix measures
interpolation between known drug--target pairs, not the generalisation a screening
campaign actually needs, and the gap between \affinityRandomCi{} and
\affinityColddrugCi{} is the size of that distinction.

\subsection{Shape: topology and curvature}
\label{sec:shape}

\begin{figure}[t]
\centering
\includegraphics[width=\linewidth]{fig11_shape_signatures}
\caption{Left: H$_1$ features of three conformers, all far above the diagonal.
Centre: Ollivier--Ricci curvature per bond; benzene is a single point on zero.
Right: long-lived H$_1$ features against perceived ring count.}
\label{fig:shape}
\end{figure}

The pipeline computes two shape descriptions that no descriptor in
Section~\ref{sec:molnet} captures, and both are checked against something that
is true independently of the code.

Persistent homology is computed on the 3D conformer and never sees the bond
graph. On \shapeProbes{} probe structures it recovers the expected number of
long-lived H$_1$ features in \shapeProbesMatched{} of \shapeProbes{} cases ---
benzene one, acyclic octane none, caffeine and naphthalene two. The interesting
test is the one that is not hand-picked: across \shapeRecoveryN{} molecules
from ESOL, the count of features surviving more than \shapeLifeThreshold{}\,Å
correlates with the ring count perceived from connectivity at
$r = \shapeRecoveryR{}$, and agrees \emph{exactly} on \shapeRecoveryExact{} of
them (mean absolute error \shapeRecoveryMae{} rings). Geometry alone recovers a
combinatorial property of the molecular graph.

Ollivier--Ricci curvature is computed on the bond graph with an exact
transportation simplex. For a cycle graph $C_n$ with $n \geq 6$ the curvature is
exactly zero, which is a theorem rather than a fit; on \shapeCycles{} ring sizes
the implementation returns $0.00 \times 10^{0}$ on every bond, in
\shapeCyclesZero{} of \shapeCycles{} cases. That exactness is what makes the
quantity readable in the panel: an unstrained ring is flat, and anything that is
not flat is strain or fusion rather than numerical drift.

\subsection{Calibrated uncertainty}
\label{sec:conformal}

\begin{figure}[t]
\centering
\includegraphics[width=\linewidth]{fig10_adaptive_conformal}
\caption{Left: mean interval width, labelled with the marginal coverage each score
achieved. Right: coverage within difficulty quintiles. Flat at nominal is ideal.}
\label{fig:adaptive}
\end{figure}

Conformal intervals hit their guarantee. At a nominal \conformalNominal{}, measured
coverage over \conformalSeeds{} seeds is \confEsolStdCov{} on ESOL,
\confLipoStdCov{} on Lipo and \confFreesolvStdCov{} on FreeSolv. The FreeSolv figure
is over-coverage, not a failure: the guarantee is one-sided.

The interesting question is width. With the default score every molecule receives an
interval of the same length, which cannot distinguish a confident prediction from a
guess. Conditional coverage shows the cost: on Lipo, coverage drifts by
\confLipoStdSpread{} across difficulty quintiles --- over-covering easy molecules and
under-covering hard ones. Scaling the nonconformity score by the Gaussian-process
posterior standard deviation cuts that drift to \confLipoSigmaSpread{} and narrows
the mean interval from \confLipoStdWidth{} to \confLipoSigmaWidth{}, at unchanged
marginal coverage (\confLipoSigmaCov{}). On ESOL the same change makes conditional
coverage \emph{worse}, from \confEsolStdSpread{} to \confEsolSigmaSpread{}. The
adaptive score is therefore offered as an option and left off by default.

\subsection{Gene Ontology enrichment, and the ordering that breaks it}
\label{sec:go}

\begin{figure*}[t]
\centering
\includegraphics[width=\linewidth]{fig13_go_enrichment_validation}
\caption{Left: terms called significant on random gene sets, which should be none.
Correcting over only the terms that happen to have hits is well calibrated on large
sets and badly miscalibrated on small ones, which is where enrichment is usually
run. Centre and right: detection and rank for a term planted into a
\goStudySize{}-gene set.}
\label{fig:go}
\end{figure*}

The panel ships a pruned ontology: biological-process terms with between
\goWindowLo{} and \goWindowHi{} annotated genes, \goTermsKept{} of \goTermsTotal{}
terms over \goGenes{} gene symbols, interned as indices to fit the package limit at
\goIndexMb{}~MB. Over-representation is hypergeometric, corrected across terms with
Benjamini--Hochberg \cite{benjamini1995controlling} at
$\mathrm{FDR}=\goAlpha{}$.

Two conditions matter, and only the second is usually reported. An enrichment must
recover biology that is there, and it must stay silent when there is none. The
second is where this implementation was wrong. Terms were first dropped for having
too few hits, and only the survivors entered the correction --- which makes the
denominator depend on the outcomes being corrected. Both orderings are run here on
identical random sets (Figure~\ref{fig:go}, left). On \goSmallSet{}-gene sets the
defective ordering returns a mean of \goSmallBad{} significant terms drawn from
nothing but noise, as many as \goSmallBadMax{} in one trial, with at least one term
in \goSmallBadTrials{} of \goTrials{} trials. Correcting across every term in the
window, including the ones that came back empty, gives \goSmallGood{}.

The defect is invisible at the sizes one is most likely to test with: by
\goLargeSet{}-gene sets the two orderings agree (\goLargeBad{} against
\goLargeGood{}). It bites hardest on short, focused lists, because those leave few
terms with any hits at all and the correction's denominator collapses onto the
handful most likely to look significant. Across every set size, the shipped
ordering never exceeds \goWorstGood{} terms per random set.

Sensitivity is measured by planting $k$ genes from one term into a
\goStudySize{}-gene set, over \goPlantedTerms{} terms (Figure~\ref{fig:go}, centre
and right). The floor is sharp: \goFloorK{} planted genes are enough for every
term to be detected and ranked inside the top ten, while \goLowK{} are recovered
only \goLowRate{} of the time and land near rank thirty, far below anything a
reader would scroll to. Diluting a whole term, rather than shrinking it, barely
registers --- terms still rank first with most of the set replaced by noise --- so
the floor, not a dilution curve, is the honest statement of what the method needs.

Applied to the kinase inhibitor map, \goTargetsResolvedNames{} of
\goTargetsTotal{} target names resolve, collapsing to \goTargetsMatched{} distinct
genes --- a panel writes one kinase under several names, so the two counts are
reported separately rather than conflated. Names are resolved by transformations
that are accepted only when they land on a real gene, so a mutation suffix
(\texttt{KIT(D816V)}), a phospho form (\texttt{ABL1p}) or a complex
(\texttt{CDK4-cyclinD1}) reaches its parent gene and nothing else can be silently
substituted. The remainder are previous HGNC symbols such as \texttt{ASK1} and
\texttt{ARK5}; they are reported as unmatched rather than guessed at, since a
mapping invented to raise the match rate is indistinguishable in the output from
biology that was measured.
